%%%
%%% Radical "⽿"
%%%
\section*{Radical 128: ``⽿''}\addcontentsline{toc}{section}{Radical 128: ⽿}\addcontentsline{loh}{figure}{\#\#\#\# 128: ⽿}

%%%%%%%%%% 耳 %%%%%%%%%%
\subsection*{耳}\addcontentsline{loh}{figure}{耳}

\begin{Entry}{耳}{6}{⽿}[Kangxi 128]
  \begin{Phonetics}{耳}{er3}
    \definition*{s.}{Sobrenome: Er}
    \definition{part.}{(clássico) somente; apenas}
    \definition{s.}{orelha | coisa parecida com uma orelha | em ambos os lados; lado | orelha de um utensílio}
  \end{Phonetics}
\end{Entry}

\begin{Entry}{耳目一新}{6,5,1,13}{⽿,⽬,⼀,⽄}
  \begin{Phonetics}{耳目一新}{er3mu4-yi4xin1}[][HSK 7-9]
    \definition{expr.}{encontrar tudo fresco e novo; encontrar"-se em um mundo inteiramente novo; apresentar uma nova aparência (de um lugar); uma mudança agradável de atmosfera; ``Tudo o que ouço e vejo mudou e parece novo.''}
  \end{Phonetics}
\end{Entry}

\begin{Entry}{耳光}{6,6}{⽿,⼉}
  \begin{Phonetics}{耳光}{er3guang1}[][HSK 7-9]
    \definition[个,记]{s.}{uma bofetada na orelha; um tapa na cara; (bater) no rosto em frente à orelha; a ação de bater no rosto}
  \end{Phonetics}
\end{Entry}

\begin{Entry}{耳朵}{6,6}{⽿,⽊}
  \begin{Phonetics}{耳朵}{er3duo5}[][HSK 5]
    \definition[双,只,个,对]{s.}{orelha; ouvido; órgão da audição e do equilíbrio}
  \end{Phonetics}
\end{Entry}

\begin{Entry}{耳机}{6,6}{⽿,⽊}
  \begin{Phonetics}{耳机}{er3ji1}[][HSK 4]
    \definition[副,个,对]{s.}{fone de ouvido; receptor (de telefone); dispositivos que permitem que uma pessoa ouça sons sozinha, como ouvir música, histórias, chamadas telefônicas etc., usados na cabeça ou inseridos nos ouvidos}
  \end{Phonetics}
\end{Entry}

\begin{Entry}{耳闻目睹}{6,9,5,13}{⽿,⾨,⽬,⽬}
  \begin{Phonetics}{耳闻目睹}{er3wen2-mu4du3}[][HSK 7-9]
    \definition{expr.}{testemunhar pessoalmente; ver e ouvir pessoalmente; o que se vê e se ouve}
  \end{Phonetics}
\end{Entry}

\begin{Entry}{耳熟能详}{6,15,10,8}{⽿,⽕,⾁,⾔}
  \begin{Phonetics}{耳熟能详}{er3shu2-neng2xiang2}[][HSK 7-9]
    \definition{expr.}{o que é ouvido com frequência pode ser repetido em detalhes; já ouvi isso muitas vezes e estou familiarizado o suficiente para falar sobre isso em detalhes}
  \end{Phonetics}
\end{Entry}

%%%%%%%%%% 耸 %%%%%%%%%%
\subsection*{耸}\addcontentsline{loh}{figure}{耸}

\begin{Entry}{耸}{10}{⽿}
  \begin{Phonetics}{耸}{song3}
    \definition{v.}{ser imponente; ser elevado | alarmar; chocar | encolher os ombros; levantar (os ombros, etc.) | elevar"-se; erguer"-se verticalmente | chocar; atrair a atenção; surpreender}
  \end{Phonetics}
\end{Entry}

\begin{Entry}{耸立}{10,5}{⽿,⽴}
  \begin{Phonetics}{耸立}{song3li4}[][HSK 7-9]
    \definition{v.}{elevar"-se no alto; erguer"-se bem alto; manter"-se ereto}
  \synonymref{操场}{cao1chang3}
  \synonymref{改为}{gai3wei2}
  \synonymref{挺拔}{ting3ba2}
  \synonymref{挺立}{ting3li4}
  \synonymref{屹立}{yi4li4}
  \synonymref{坐落}{zuo4luo4}
  \end{Phonetics}
\end{Entry}

%%%%%%%%%% 耻 %%%%%%%%%%
\subsection*{耻}\addcontentsline{loh}{figure}{耻}

\begin{Entry}{耻}{10}{⽿}
  \begin{Phonetics}{耻}{chi3}
    \definition{s.}{vergonha; desgraça; humilhação}
    \definition{v.}{estar envergonhado de; considerar vergonhoso}
  \end{Phonetics}
\end{Entry}

\begin{Entry}{耻笑}{10,10}{⽿,⽵}
  \begin{Phonetics}{耻笑}{chi3xiao4}[][HSK 7-9]
    \definition{v.}{ridicularizar alguém; zombar; zombar de; rir de}
  \end{Phonetics}
\end{Entry}

\begin{Entry}{耻辱}{10,10}{⽿,⾠}
  \begin{Phonetics}{耻辱}{chi3ru3}[][HSK 7-9]
    \definition{s.}{vergonha; desgraça; humilhação; danos à reputação; incidente vergonhoso}
  \end{Phonetics}
\end{Entry}

%%%%%%%%%% 耽 %%%%%%%%%%
\subsection*{耽}\addcontentsline{loh}{figure}{耽}

\begin{Entry}{耽}{10}{⽿}
  \begin{Phonetics}{耽}{dan1}
    \definition*{s.}{Sobrenome: Dan}
    \definition{v.}{atrasar | Literário: abandonar"-se a; entregar"-se a}
  \end{Phonetics}
\end{Entry}

\begin{Entry}{耽心}{10,4}{⽿,⼼}
  \begin{Phonetics}{耽心}{dan1xin1}
    \variantof{担心}
  \end{Phonetics}
\end{Entry}

\begin{Entry}{耽误}{10,9}{⽿,⾔}
  \begin{Phonetics}{耽误}{dan1wu5}[][HSK 7-9]
    \definition{v.}{atrasar; segurar; perder algo devido a atraso ou oportunidade perdida; perder (oportunidade)}
  \end{Phonetics}
\end{Entry}

\begin{Entry}{耽搁}{10,12}{⽿,⼿}
  \begin{Phonetics}{耽搁}{dan1ge5}[][HSK 7-9]
    \definition{v.}{ficar; fazer uma parada | atrasar | perder (uma oportunidade, um prazo)}
  \end{Phonetics}
\end{Entry}

%%%%%%%%%% 耿 %%%%%%%%%%
\subsection*{耿}\addcontentsline{loh}{figure}{耿}

\begin{Entry}{耿}{10}{⽿}
  \begin{Phonetics}{耿}{geng3}
    \definition*{s.}{Sobrenome: Geng}
    \definition{adj.}{Literário: brilhante | honesto e justo; correto; íntegro | dedicado; leal}
  \end{Phonetics}
\end{Entry}

\begin{Entry}{耿直}{10,8}{⽿,⽬}
  \begin{Phonetics}{耿直}{geng3zhi2}[][HSK 7-9]
    \definition{adj.}{íntregro; franco; correto; honesto e franco}
  \end{Phonetics}
\end{Entry}

%%%%%%%%%% 聊 %%%%%%%%%%
\subsection*{聊}\addcontentsline{loh}{figure}{聊}

\begin{Entry}{聊}{11}{⽿}
  \begin{Phonetics}{聊}{liao2}[][HSK 6]
    \definition*{s.}{Sobrenome: Liao}
    \definition{adv.}{apenas; meramente; provisoriamente; por enquanto | um pouco; ligeiramente}
    \definition{v.}{tagarelar; conversar; bater papo | confiar (ou depender, recorrer) a}
  \synonymref{粗}{cu1}
  \synonymref{靠}{kao4}
  \synonymref{赖}{lai4}
  \synonymref{略}{lve4}
  \synonymref{且}{qie3}
  \synonymref{稍}{shao1}
  \synonymref{谈}{tan2}
  \synonymref{依}{yi1}
  \end{Phonetics}
\end{Entry}

\begin{Entry}{聊天}{11,4}{⽿,⼤}
  \begin{Phonetics}{聊天}{liao2//tian1}
    \definition{v.+compl.}{papear; bater papo; fofocar}
  \end{Phonetics}
\end{Entry}

\begin{Entry}{聊天儿}{11,4,2}{⽿,⼤,⼉}
  \begin{Phonetics}{聊天儿}{liao2//tian1r5}[][HSK 6]
    \definition{v.+compl.}{conversar; fofocar; bater papo; duas ou mais pessoas conversando sem um tópico ou propósito específico}
  \synonymref{交谈}{jiao1tan2}
  \synonymref{聊}{liao2}
  \end{Phonetics}
\end{Entry}

%%%%%%%%%% 聋 %%%%%%%%%%
\subsection*{聋}\addcontentsline{loh}{figure}{聋}

\begin{Entry}{聋}{11}{⽿}
  \begin{Phonetics}{聋}{long2}[][HSK 7-9]
    \definition{adj.}{surdo; com deficiência auditiva}
  \end{Phonetics}
\end{Entry}

\begin{Entry}{聋人}{11,2}{⽿,⼈}
  \begin{Phonetics}{聋人}{long2ren2}[][HSK 7-9]
    \definition[个,位,名]{s.}{surdo (pessoa surda)}
  \end{Phonetics}
\end{Entry}

%%%%%%%%%% 职 %%%%%%%%%%
\subsection*{职}\addcontentsline{loh}{figure}{职}

\begin{Entry}{职}{11}{⽿}
  \begin{Phonetics}{职}{zhi2}
    \definition*{s.}{Sobrenome: Zhi}
    \definition{prep.}{para; devido a; por causa de}
    \definition{prep.}{Obsoleto: Eu (em relatórios oficiais aos superiores)}
    \definition{s.}{dever; trabalho | cargo; posto; função; responsabilidades; posição}
    \definition{v.}{gerenciar; dirigir | administrar}
  \end{Phonetics}
\end{Entry}

\begin{Entry}{职工}{11,3}{⽿,⼯}
  \begin{Phonetics}{职工}{zhi2gong1}[][HSK 3]
    \definition[个,位,名,些]{s.}{pessoal; trabalhadores e funcionários administrativos}
  \synonymref{工人}{gong1ren5}
  \synonymref{员工}{yuan2gong1}
  \synonymref{职员}{zhi2yuan2}
  \end{Phonetics}
\end{Entry}

\begin{Entry}{职业}{11,5}{⽿,⼀}
  \begin{Phonetics}{职业}{zhi2ye4}[][HSK 3]
    \definition{adj.}{profissional; não amador}
    \definition[种,份,个]{s.}{ocupação; profissão; vocação; o trabalho que um indivíduo realiza na sociedade como sua principal fonte de subsistência}
  \synonymref{工作}{gong1zuo4}
  \synonymref{行业}{hang2ye4}
  \synonymref{事业}{shi4ye4}
  \synonymref{职务}{zhi2wu4}
  \synonymref{专长}{zhuan1chang2}
  \synonymref{专业}{zhuan1ye4}
  \synonymref{专职}{zhuan1zhi2}
  \antonymref{消遣}{xiao1qian3}
  \antonymref{休闲}{xiu1xian2}
  \antonymref{业余}{ye4yu2}
  \end{Phonetics}
\end{Entry}

\begin{Entry}{职业病}{11,5,10}{⽿,⼀,⽧}
  \begin{Phonetics}{职业病}{zhi2ye4bing4}[][HSK 7-9]
    \definition[个]{s.}{doença ocupacional; doenças crônicas causadas pela natureza de certos trabalhos ou ambientes de trabalho específicos, como a pneumoconiose, comum entre mineiros e trabalhadores da indústria cerâmica, e o enfisema, comum entre sopradores de vidro}
  \end{Phonetics}
\end{Entry}

\begin{Entry}{职务}{11,5}{⽿,⼒}
  \begin{Phonetics}{职务}{zhi2wu4}[][HSK 5]
    \definition{s.}{cargo; posto; deveres; função; funções que devem ser desempenhadas de acordo com as especificações do cargo}
  \synonymref{岗位}{gang3wei4}
  \synonymref{职位}{zhi2wei4}
  \synonymref{职责}{zhi2ze2}
  \end{Phonetics}
\end{Entry}

\begin{Entry}{职权}{11,6}{⽿,⽊}
  \begin{Phonetics}{职权}{zhi2quan2}[][HSK 7-9]
    \definition[项]{s.}{função e poder; poderes do cargo; autoridade do cargo}
  \synonymref{权利}{quan2li4}
  \synonymref{权力}{quan2li4}
  \end{Phonetics}
\end{Entry}

\begin{Entry}{职位}{11,7}{⽿,⼈}
  \begin{Phonetics}{职位}{zhi2wei4}[][HSK 5]
    \definition[个]{s.}{posto; posição; cargo que exerce determinadas funções em órgãos ou entidades}
  \synonymref{地位}{di4wei4}
  \synonymref{位置}{wei4zhi4}
  \synonymref{位子}{wei4zi5}
  \synonymref{职务}{zhi2wu4}
  \end{Phonetics}
\end{Entry}

\begin{Entry}{职员}{11,7}{⽿,⼝}
  \begin{Phonetics}{职员}{zhi2yuan2}[][HSK 7-9]
    \definition[个,位,名,些]{s.}{funcionário; trabalhador de escritório; membro da equipe; pessoas que ocupam cargos administrativos ou operacionais em agências governamentais, empresas, escolas e organizações}
  \seealsoref{工人}{gong1ren5}
  \synonymref{人员}{ren2yuan2}
  \synonymref{职工}{zhi2gong1}
  \end{Phonetics}
\end{Entry}

\begin{Entry}{职责}{11,8}{⽿,⾙}
  \begin{Phonetics}{职责}{zhi2ze2}[][HSK 6]
    \definition[种]{s.}{dever; obrigação; responsabilidade; coisas que você deve fazer por causa de sua profissão ou identidade}
  \synonymref{本分}{ben3fen4}
  \synonymref{工作}{gong1zuo4}
  \synonymref{任务}{ren4wu4}
  \synonymref{使命}{shi3ming4}
  \synonymref{义务}{yi4wu4}
  \synonymref{责任}{ze2ren4}
  \synonymref{职务}{zhi2wu4}
  \end{Phonetics}
\end{Entry}

\begin{Entry}{职能}{11,10}{⽿,⾁}
  \begin{Phonetics}{职能}{zhi2neng2}[][HSK 5]
    \definition[种,项]{s.}{função; funções ou papéis que as organizações, instituições, etc. devem desempenhar}
  \synonymref{功能}{gong1neng2}
  \synonymref{性能}{xing4neng2}
  \synonymref{职责}{zhi2ze2}
  \synonymref{作用}{zuo4yong4}
  \end{Phonetics}
\end{Entry}

%%%%%%%%%% 联 %%%%%%%%%%
\subsection*{联}\addcontentsline{loh}{figure}{联}

\begin{Entry}{联}{12}{⽿}
  \begin{Phonetics}{联}{lian2}
    \definition{s.}{dísticos (antitéticos)}
    \definition{v.}{aliar"-se a; unir"-se; juntar"-se a}
  \end{Phonetics}
\end{Entry}

\begin{Entry}{联手}{12,4}{⽿,⼿}
  \begin{Phonetics}{联手}{lian2shou3}[][HSK 6]
    \definition{v.}{dar as mãos; cooperar | Literário: dar as mãos | agir em conjunto}
  \synonymref{合伙}{he2huo3}
  \synonymref{合作}{he2zuo4}
  \synonymref{联合}{lian2he2}
  \synonymref{协同}{xie2tong2}
  \synonymref{协作}{xie2zuo4}
  \antonymref{对立}{dui4li4}
  \antonymref{分裂}{fen1lie4}
  \end{Phonetics}
\end{Entry}

\begin{Entry}{联合}{12,6}{⽿,⼝}
  \begin{Phonetics}{联合}{lian2he2}[][HSK 3]
    \definition{adj.}{conjunto; unido; federal; combinado}
    \definition{s.}{aliado; união; aliança; conectar"-se ou unir"-se para agir em conjunto}
  \synonymref{共同}{gong4tong2}
  \synonymref{合伙}{he2huo3}
  \synonymref{结合}{jie2he2}
  \synonymref{连接}{lian2jie1}
  \synonymref{联络}{lian2luo4}
  \synonymref{统一}{tong3yi1}
  \synonymref{团结}{tuan2jie2}
  \synonymref{协同}{xie2tong2}
  \synonymref{整合}{zheng3he2}
  \antonymref{单独}{dan1du2}
  \antonymref{对峙}{dui4zhi4}
  \antonymref{分开}{fen1//kai1}
  \antonymref{分裂}{fen1lie4}
  \antonymref{分散}{fen1san4}
  \antonymref{孤立}{gu1li4}
  \antonymref{解体}{jie3ti3}
  \antonymref{排斥}{pai2chi4}
  \end{Phonetics}
\end{Entry}

\begin{Entry}{联合会}{12,6,6}{⽿,⼝,⼈}
  \begin{Phonetics}{联合会}{lian2he2hui4}
    \definition{s.}{federação; união; associação; liga}
  \end{Phonetics}
\end{Entry}

\begin{Entry}{联合国}{12,6,8}{⽿,⼝,⼞}
  \begin{Phonetics}{联合国}{lian2he2guo2}[][HSK 3]
    \definition*{s.}{Nações Unidas; Organização internacional fundada em 1945, após o fim da Segunda Guerra Mundial, com sede em Nova Iorque, Estados Unidos ; as suas principais instituições são a Assembleia Geral, o Conselho de Segurança, o Conselho Econômico e Social e o Secretariado; de acordo com a Carta das Nações Unidas, os seus principais objetivos são manter a paz e a segurança internacionais, desenvolver relações amigáveis entre os países e promover a cooperação internacional nas áreas econômica e cultural}
  \end{Phonetics}
\end{Entry}

\begin{Entry}{联欢}{12,6}{⽿,⽋}
  \begin{Phonetics}{联欢}{lian2huan1}[][HSK 7-9]
    \definition{v.}{ter um encontro; fazer uma reunião; (membros de um grupo ou de dois ou mais grupos) reunir"-se para celebrar ou fortalecer a união}
  \end{Phonetics}
\end{Entry}

\begin{Entry}{联网}{12,6}{⽿,⽹}
  \begin{Phonetics}{联网}{lian2//wang3}[][HSK 7-9]
    \definition{pref.}{cyber-}
    \definition{v.+compl.}{conectar computadores (ou fios elétricos, linhas de comunicação, etc.) para formar uma rede; estar \emph{on"-line}; estar conectado}
  \end{Phonetics}
\end{Entry}

\begin{Entry}{联邦}{12,6}{⽿,⾢}
  \begin{Phonetics}{联邦}{lian2bang1}[][HSK 7-9]
    \definition{s.}{federação; união; comunidade}
  \end{Phonetics}
\end{Entry}

\begin{Entry}{联系}{12,7}{⽿,⽷}
  \begin{Phonetics}{联系}{lian2xi4}[][HSK 3]
    \definition[个,种,层]{s.}{relacionamento; relacionamento entre duas coisas}
    \definition{v.}{entrar em contato; contatar; comunicar"-se com alguém por telefone, \emph{e"-mail} ou carta | agendar; entrar em contato com; estabelecer algum tipo de relação com a outra parte | relacionar; combinar; integrar}
  \seealsoref{联络}{lian2luo4}
  \synonymref{沟通}{gou1tong1}
  \synonymref{挂钩}{gua4gou1}
  \synonymref{关联}{guan1lian2}
  \synonymref{关系}{guan1xi5}
  \synonymref{纽带}{niu3dai4}
  \synonymref{相关}{xiang1guan1}
  \antonymref{分割}{fen1ge1}
  \antonymref{孤立}{gu1li4}
  \antonymref{脱节}{tuo1jie2}
  \antonymref{脱离}{tuo1li2}
  \end{Phonetics}
\end{Entry}

\begin{Entry}{联络}{12,9}{⽿,⽷}
  \begin{Phonetics}{联络}{lian2luo4}[][HSK 5]
    \definition{v.}{entrar em contato; comunicar"-se; entrar em contato com}
  \seealsoref{联系}{lian2xi4}
  \synonymref{拉拢}{la1long3}
  \synonymref{联合}{lian2he2}
  \synonymref{连接}{lian2jie1}
  \synonymref{团结}{tuan2jie2}
  \antonymref{切断}{qie1duan4}
  \antonymref{中断}{zhong1duan4}
  \end{Phonetics}
\end{Entry}

\begin{Entry}{联想}{12,13}{⽿,⼼}
  \begin{Phonetics}{联想}{lian2xiang3}[][HSK 5]
    \definition*{s.}{Lenovo (empresa)}
    \definition{v.}{associar"-se a; estabelecer uma conexão mental; lembrar"-se de algo; lembrar"-se de outras pessoas ou coisas relacionadas devido a alguém ou algo; evocar outros conceitos relacionados devido a um determinado conceito}
  \synonymref{设想}{she4xiang3}
  \synonymref{想象}{xiang3xiang4}
  \end{Phonetics}
\end{Entry}

\begin{Entry}{联盟}{12,13}{⽿,⽫}
  \begin{Phonetics}{联盟}{lian2meng2}[][HSK 6]
    \definition{s.}{aliança; coalizão; liga; união}
  \synonymref{联邦}{lian2bang1}
  \synonymref{同盟}{tong2meng2}
  \antonymref{分裂}{fen1lie4}
  \end{Phonetics}
\end{Entry}

\begin{Entry}{联赛}{12,14}{⽿,⾙}
  \begin{Phonetics}{联赛}{lian2sai4}[][HSK 6]
    \definition{s.}{jogos da liga | liga (esportiva) | torneio da liga}
  \end{Phonetics}
\end{Entry}

%%%%%%%%%% 聘 %%%%%%%%%%
\subsection*{聘}\addcontentsline{loh}{figure}{聘}

\begin{Entry}{聘}{13}{⽿}
  \begin{Phonetics}{聘}{pin4}[][HSK 7-9]
    \definition{v.}{contratar | visitar como enviado; visitar como representante da família para fins de arranjo matrimonial (cultura tradicional) | (uma menina) casar"-se ou ser dada em casamento; ficar noivos}
  \end{Phonetics}
\end{Entry}

\begin{Entry}{聘用}{13,5}{⽿,⽤}
  \begin{Phonetics}{聘用}{pin4yong4}[][HSK 7-9]
    \definition{v.}{contratar; nomear alguém para um cargo}
  \end{Phonetics}
\end{Entry}

\begin{Entry}{聘任}{13,6}{⽿,⼈}
  \begin{Phonetics}{聘任}{pin4ren4}[][HSK 7-9]
    \definition{v.}{contratar alguém para; nomear alguém para um cargo; contratar; recrutar por convite; contratar para atuar como (cargo)}
  \end{Phonetics}
\end{Entry}

\begin{Entry}{聘请}{13,10}{⽿,⾔}
  \begin{Phonetics}{聘请}{pin4qing3}[][HSK 6]
    \definition{v.}{convidar; empregar; envolver; chamar; contratar alguém para assumir uma posição}
  \synonymref{雇佣}{gu4yong1}
  \synonymref{聘任}{pin4ren4}
  \synonymref{聘用}{pin4yong4}
  \synonymref{邀请}{yao1qing3}
  \synonymref{招聘}{zhao1pin4}
  \antonymref{辞退}{ci2tui4}
  \antonymref{解雇}{jie3gu4}
  \end{Phonetics}
\end{Entry}

%%%%%%%%%% 聚 %%%%%%%%%%
\subsection*{聚}\addcontentsline{loh}{figure}{聚}

\begin{Entry}{聚}{14}{⽿}
  \begin{Phonetics}{聚}{ju4}[][HSK 4]
    \definition*{s.}{Sobrenome: Ju}
    \definition{v.}{reunir"-se; juntar"-se}
  \seealsoref{聚会}{ju4hui4}
  \synonymref{合}{he2}
  \synonymref{汇}{hui4}
  \synonymref{积}{ji1}
  \synonymref{集}{ji2}
  \antonymref{分}{fen1}
  \antonymref{散}{san4}
  \end{Phonetics}
\end{Entry}

\begin{Entry}{聚会}{14,6}{⽿,⼈}
  \begin{Phonetics}{聚会}{ju4hui4}[][HSK 4]
    \definition[个,次]{s.}{reunião; encontro; confraternização; festa}
    \definition{v.}{encontrar"-se; reunir"-se}
  \seealsoref{聚}{ju4}
  \synonymref{集合}{ji2he2}
  \synonymref{集会}{ji2hui4}
  \synonymref{集结}{ji2jie2}
  \synonymref{集中}{ji2zhong1}
  \synonymref{聚集}{ju4ji2}
  \synonymref{开会}{kai1//hui4}
  \synonymref{团聚}{tuan2ju4}
  \synonymref{相聚}{xiang1ju4}
  \antonymref{分离}{fen1li2}
  \antonymref{分散}{fen1san4}
  \antonymref{解散}{jie3san4}
  \end{Phonetics}
\end{Entry}

\begin{Entry}{聚散}{14,12}{⽿,⽁}
  \begin{Phonetics}{聚散}{ju4san4}
    \definition{s.}{agregação e dissipação}
    \definition{v.}{encontrar e separar | unir"-se e separar"-se}
  \end{Phonetics}
\end{Entry}

\begin{Entry}{聚集}{14,12}{⽿,⾫}
  \begin{Phonetics}{聚集}{ju4ji2}[][HSK 7-9]
    \definition{v.}{reunir; juntar; coletar; reunir"-se; juntar"-se}
  \end{Phonetics}
\end{Entry}

\begin{Entry}{聚精会神}{14,14,6,9}{⽿,⽶,⼈,⽰}
  \begin{Phonetics}{聚精会神}{ju4jing1-hui4shen2}[][HSK 7-9]
    \definition{expr.}{concentrado; concentrar a atenção; focar a mente; estar absorto em; estar profundamente concentrado; estar totalmente concentrado}
  \end{Phonetics}
\end{Entry}

%%%%%%%%%% 聪 %%%%%%%%%%
\subsection*{聪}\addcontentsline{loh}{figure}{聪}

\begin{Entry}{聪}{15}{⽿}
  \begin{Phonetics}{聪}{cong1}
    \definition{adj.}{audição aguçada | brilhante; inteligente; esperto | perspicaz}
    \definition{s.}{(literário) faculdades auditivas}
  \end{Phonetics}
\end{Entry}

\begin{Entry}{聪明}{15,8}{⽿,⽇}
  \begin{Phonetics}{聪明}{cong1ming5}[][HSK 5]
    \definition{adj.}{brilhante; esperto; inteligente; intelecto bem desenvolvido com boa memória e capacidade de compreensão}
  \seealsoref{机灵}{ji1ling5}
  \seealsoref{伶俐}{ling2li4}
  \synonymref{聪慧}{cong1hui4}
  \synonymref{机智}{ji1zhi4}
  \synonymref{精明}{jing1ming2}
  \synonymref{灵活}{ling2huo2}
  \synonymref{灵敏}{ling2min3}
  \synonymref{明智}{ming2zhi4}
  \synonymref{智慧}{zhi4hui4}
  \antonymref{笨蛋}{ben4dan4}
  \antonymref{笨拙}{ben4zhuo1}
  \antonymref{糊涂}{hu2tu5}
  \antonymref{傻瓜}{sha3gua1}
  \antonymref{愚蠢}{yu2chun3}
  \antonymref{愚昧}{yu2mei4}
  \end{Phonetics}
\end{Entry}

\begin{Entry}{聪慧}{15,15}{⽿,⼼}
  \begin{Phonetics}{聪慧}{cong1hui4}
    \definition{adj.}{brilhante; inteligente; sábio}
  \synonymref{聪明}{cong1ming5}
  \synonymref{机灵}{ji1ling5}
  \synonymref{伶俐}{ling2li4}
  \synonymref{灵敏}{ling2min3}
  \synonymref{灵巧}{ling2qiao3}
  \synonymref{智慧}{zhi4hui4}
  \antonymref{痴呆}{chi1dai1}
  \antonymref{愚蠢}{yu2chun3}
  \end{Phonetics}
\end{Entry}

%%%%% EOF %%%%%

